% GENERATED FILE. Edit canonical lessons, then run npm run generate:books.
% canonical-source-hash: fnv1a64:e34eea6c753e8582
% canonical-lessons: TA-C248-varatci, TA-C248-nilanatukkam, TA-C248-iyarkai, TA-C248-parai, TA-C248-neruppu, TA-R248-first-pass-words-from-chapters-243-245, TA-R248-second-pass-words-from-chapters-246-248

\chapter{Drought, Earthquake, Nature, Rock, Fire}
\label{ch:ta-drought-earthquake-nature-rock-fire}
\hlchaptermodality{248}

\begin{chapteropening}
\textbf{By the end of this chapter:} \emph{I can say a drought, an earthquake, nature, a rock and fire.}

Complete the last lesson of chapter 248: I can say a drought, an earthquake, nature, a rock and fire.
\end{chapteropening}

\section[\texorpdfstring{\ta{வறட்சி} (vaṟaṭci)}{வறட்சி (vaṟaṭci)}]{\ta{வறட்சி} (vaṟaṭci) — a drought}

\label{lesson:TA-C248-varatci}



\begin{quote}
Before the new one: say the Tamil for a squirrel, then the Tamil for a worm.
\end{quote}

\subsection*{You'll want to know: \ta{வறட்சி}}
\textbf{\ta{வறட்சி}} — \emph{vaṟaṭci} — \textquotedblleft{}a drought\textquotedblright{}.

\begin{grammarlens}[title={What you've built}]
One more word for this chapter.
\end{grammarlens}

\subsection*{Your turn}
\begin{itemize}
\raggedright
  \item \emph{Say it:} \emph{vaṟaṭci}
  \item \emph{Say it:} \emph{vaṟaṭci}, once more
  \item \emph{Say it:} say \emph{vaṟaṭci}
  \item \emph{Recall:} say the Tamil for a squirrel, then the Tamil for a worm, then say \emph{vaṟaṭci} again
\end{itemize}

\begin{tcolorbox}[breakable,colback=teal!4,colframe=teal!35!black,title={Before you move on}]
What does \ta{வறட்சி} mean? (\textquotedblleft{}A drought\textquotedblright{}.) Say it once more.
\end{tcolorbox}
\section[\texorpdfstring{\ta{நிலநடுக்கம்} (nilanaṭukkam)}{நிலநடுக்கம் (nilanaṭukkam)}]{\ta{நிலநடுக்கம்} (nilanaṭukkam) — an earthquake}

\label{lesson:TA-C248-nilanatukkam}



\begin{quote}
Before the new one: say the Tamil for a worm, then the Tamil for a drought.
\end{quote}

\subsection*{You'll want to know: \ta{நிலநடுக்கம்}}
\textbf{\ta{நிலநடுக்கம்}} — \emph{nilanaṭukkam} — \textquotedblleft{}an earthquake\textquotedblright{}.

\begin{grammarlens}[title={What you've built}]
One more word for this chapter.
\end{grammarlens}

\subsection*{Your turn}
\begin{itemize}
\raggedright
  \item \emph{Say it:} \emph{nilanaṭukkam}
  \item \emph{Say it:} \emph{nilanaṭukkam}, once more
  \item \emph{Say it:} say \emph{nilanaṭukkam}
  \item \emph{Recall:} say the Tamil for a worm, then the Tamil for a drought, then say \emph{nilanaṭukkam} again
\end{itemize}

\begin{tcolorbox}[breakable,colback=teal!4,colframe=teal!35!black,title={Before you move on}]
What does \ta{நிலநடுக்கம்} mean? (\textquotedblleft{}An earthquake\textquotedblright{}.) Say it once more.
\end{tcolorbox}
\section[\texorpdfstring{\ta{இயற்கை} (iyaṟkai)}{இயற்கை (iyaṟkai)}]{\ta{இயற்கை} (iyaṟkai) — nature}

\label{lesson:TA-C248-iyarkai}



\begin{quote}
Before the new one: say the Tamil for a drought, then the Tamil for an earthquake.
\end{quote}

\subsection*{You'll want to know: \ta{இயற்கை}}
\textbf{\ta{இயற்கை}} — \emph{iyaṟkai} — \textquotedblleft{}nature\textquotedblright{}.

\begin{grammarlens}[title={What you've built}]
One more word for this chapter.
\end{grammarlens}

\subsection*{Your turn}
\begin{itemize}
\raggedright
  \item \emph{Say it:} \emph{iyaṟkai}
  \item \emph{Say it:} \emph{iyaṟkai}, once more
  \item \emph{Say it:} say \emph{iyaṟkai}
  \item \emph{Recall:} say the Tamil for a drought, then the Tamil for an earthquake, then say \emph{iyaṟkai} again
\end{itemize}

\begin{tcolorbox}[breakable,colback=teal!4,colframe=teal!35!black,title={Before you move on}]
What does \ta{இயற்கை} mean? (\textquotedblleft{}Nature\textquotedblright{}.) Say it once more.
\end{tcolorbox}
\section[\texorpdfstring{\ta{பாறை} (pāṟai)}{பாறை (pāṟai)}]{\ta{பாறை} (pāṟai) — a rock}

\label{lesson:TA-C248-parai}



\begin{quote}
Before the new one: say the Tamil for an earthquake, then the Tamil for nature.
\end{quote}

\subsection*{You'll want to know: \ta{பாறை}}
\textbf{\ta{பாறை}} — \emph{pāṟai} — \textquotedblleft{}a rock\textquotedblright{}.

\begin{grammarlens}[title={What you've built}]
One more word for this chapter.
\end{grammarlens}

\subsection*{Your turn}
\begin{itemize}
\raggedright
  \item \emph{Say it:} \emph{pāṟai}
  \item \emph{Say it:} \emph{pāṟai}, once more
  \item \emph{Say it:} say \emph{pāṟai}
  \item \emph{Recall:} say the Tamil for an earthquake, then the Tamil for nature, then say \emph{pāṟai} again
\end{itemize}

\begin{tcolorbox}[breakable,colback=teal!4,colframe=teal!35!black,title={Before you move on}]
What does \ta{பாறை} mean? (\textquotedblleft{}A rock\textquotedblright{}.) Say it once more.
\end{tcolorbox}
\section[\texorpdfstring{\ta{நெருப்பு} (neruppu)}{நெருப்பு (neruppu)}]{\ta{நெருப்பு} (neruppu) — fire}

\label{lesson:TA-C248-neruppu}



\begin{quote}
Before the new one: say the Tamil for nature, then the Tamil for a rock.
\end{quote}

\subsection*{You'll want to know: \ta{நெருப்பு}}
\textbf{\ta{நெருப்பு}} — \emph{neruppu} — \textquotedblleft{}fire\textquotedblright{}.

\begin{grammarlens}[title={What you've built}]
One more word for this chapter.
\end{grammarlens}

\subsection*{Your turn}
\begin{itemize}
\raggedright
  \item \emph{Say it:} \emph{neruppu}
  \item \emph{Say it:} \emph{neruppu}, once more
  \item \emph{Say it:} say \emph{neruppu}
  \item \emph{Recall:} say the Tamil for nature, then the Tamil for a rock, then say \emph{neruppu} again
\end{itemize}

\begin{tcolorbox}[breakable,colback=teal!4,colframe=teal!35!black,title={Before you move on}]
What does \ta{நெருப்பு} mean? (\textquotedblleft{}Fire\textquotedblright{}.) Say it once more.
\end{tcolorbox}
\section[(dialogue)]{Words from chapters 243-245 — a first pass}

\label{lesson:TA-R248-first-pass-words-from-chapters-243-245}



\begin{quote}
Say the words for a rock and fire.
\end{quote}

\subsection*{Your turn}
Say these aloud:

\begin{itemize}
\raggedright
  \item wheat, a potato, an aubergine, a fragrance, a cucumber
  \item a postage stamp, a parcel, a form, a document, a passport
  \item a load, a voice, a conversation, a man, a woman
\end{itemize}

\begin{tcolorbox}[breakable,colback=teal!4,colframe=teal!35!black,title={Before you move on}]
Wheat? (\textbf{\ta{கோதுமை}}, \emph{kōtumai}.) A potato? (\textbf{\ta{உருளைக்கிழங்கு}}, \emph{uruḷaikkiḻaṅku}.) An aubergine? (\textbf{\ta{கத்தரிக்காய்}}, \emph{kattarikkāy}.) A fragrance? (\textbf{\ta{வாசனை}}, \emph{vācaṉai}.) A cucumber? (\textbf{\ta{வெள்ளரிக்காய்}}, \emph{veḷḷarikkāy}.) A postage stamp? (\textbf{\ta{அஞ்சல்} \ta{தலை}}, \emph{añcal talai}.) A parcel? (\textbf{\ta{பொட்டலம்}}, \emph{poṭṭalam}.) A form? (\textbf{\ta{படிவம்}}, \emph{paṭivam}.) A document? (\textbf{\ta{ஆவணம்}}, \emph{āvaṇam}.) A passport? (\textbf{\ta{கடவுச்சீட்டு}}, \emph{kaṭavuccīṭṭu}.) A load? (\textbf{\ta{சுமை}}, \emph{cumai}.) A voice? (\textbf{\ta{குரல்}}, \emph{kural}.) A conversation? (\textbf{\ta{உரையாடல்}}, \emph{uraiyāṭal}.) A man? (\textbf{\ta{ஆண்}}, \emph{āṇ}.) A woman? (\textbf{\ta{பெண்}}, \emph{peṇ}.)
\end{tcolorbox}
\section[(dialogue)]{Words from chapters 246-248 — a second pass}

\label{lesson:TA-R248-second-pass-words-from-chapters-246-248}



\begin{quote}
Say the words for celebration and firecrackers.
\end{quote}

\subsection*{Your turn}
Say these aloud:

\begin{itemize}
\raggedright
  \item celebration, firecrackers, culture, a habit, tradition
  \item a sparrow, an eagle, a puppy, a squirrel, a worm
  \item a drought, an earthquake, nature, a rock, fire
\end{itemize}

\begin{tcolorbox}[breakable,colback=teal!4,colframe=teal!35!black,title={Before you move on}]
Celebration? (\textbf{\ta{கொண்டாட்டம்}}, \emph{koṇṭāṭṭam}.) Firecrackers? (\textbf{\ta{பட்டாசு}}, \emph{paṭṭācu}.) Culture? (\textbf{\ta{கலாச்சாரம்}}, \emph{kalāccāram}.) A habit? (\textbf{\ta{பழக்கம்}}, \emph{paḻakkam}.) Tradition? (\textbf{\ta{பாரம்பரியம்}}, \emph{pārampariyam}.) A sparrow? (\textbf{\ta{குருவி}}, \emph{kuruvi}.) An eagle? (\textbf{\ta{கழுகு}}, \emph{kaḻuku}.) A puppy? (\textbf{\ta{நாய்க்குட்டி}}, \emph{nāykkuṭṭi}.) A squirrel? (\textbf{\ta{அணில்}}, \emph{aṇil}.) A worm? (\textbf{\ta{புழு}}, \emph{puḻu}.) A drought? (\textbf{\ta{வறட்சி}}, \emph{vaṟaṭci}.) An earthquake? (\textbf{\ta{நிலநடுக்கம்}}, \emph{nilanaṭukkam}.) Nature? (\textbf{\ta{இயற்கை}}, \emph{iyaṟkai}.) A rock? (\textbf{\ta{பாறை}}, \emph{pāṟai}.) Fire? (\textbf{\ta{நெருப்பு}}, \emph{neruppu}.)
\end{tcolorbox}
